%% ============================================================
%% shared/global_infrastructure.tex
%% AWS Global Infrastructure — SAA + SAP
%% ============================================================

\servicesection{AWS Global Infrastructure}{\saatag}{Regions, AZs, Edge locations}

\begin{keyfacts}
\begin{itemize}
  \item \textbf{Region} — geographically isolated area; contains 2+ AZs (typically 3).
        Choose region based on: latency, data residency, service availability, pricing.
  \item \textbf{Availability Zone (AZ)} — one or more discrete data centres with independent
        power, cooling, networking. AZs in a region are interconnected via high-bandwidth,
        low-latency private fibre.
  \item \textbf{Edge Location / Point of Presence (PoP)} — used by CloudFront (CDN) and
        Route~53 for caching and DNS resolution closer to end users. More locations than Regions.
  \item \textbf{Local Zone} — extension of an AWS Region to metropolitan areas; for ultra-low
        latency workloads (gaming, media, real-time).
  \item \textbf{Wavelength Zone} — AWS infrastructure embedded in 5G carrier networks for
        single-digit millisecond latency to mobile devices.
  \item \textbf{Outposts} — AWS-managed rack/server hardware in your own data centre; extends
        AWS APIs on-premises.
\end{itemize}
\end{keyfacts}

\subsection{Resilience by Design}

Design for failures at every level:
\begin{itemize}
  \item \textbf{AZ-resilient} — deploy across $\geq 2$ AZs; use Multi-AZ RDS, ALB, ASG spread.
  \item \textbf{Region-resilient} — cross-region replication (S3, Aurora Global, DynamoDB Global Tables).
  \item \textbf{Global} — Route~53 latency/health-check routing, CloudFront origin failover.
\end{itemize}

\begin{examtip}
Questions about ``highest availability'' or ``survive an AZ failure'' always point to
\textbf{Multi-AZ deployments}. Questions about ``survive a region failure'' require cross-region
architectures — Aurora Global Database, S3 Cross-Region Replication, or Route~53 failover.
\end{examtip}

\subsection{Global Infrastructure Diagram}

\begin{center}
\begin{tikzpicture}[
  region/.style={draw=awsdark, fill=awsdark!8, rounded corners=6pt,
                 minimum width=5.2cm, minimum height=2.8cm, font=\small},
  az/.style={draw=saacolor, fill=saacolor!10, rounded corners=3pt,
             minimum width=1.4cm, minimum height=1.2cm, font=\tiny\bfseries, align=center},
  edge/.style={draw=tipcolor, fill=tipcolor!10, rounded corners=3pt,
               minimum width=1.6cm, minimum height=0.6cm, font=\tiny, align=center},
  lbl/.style={font=\scriptsize\bfseries, color=awsdark}
]

%% Region box
\node[region] (region) at (0,0) {};
\node[lbl] at (0, 1.2) {AWS Region};

%% Three AZs inside
\node[az] (az1) at (-1.6, 0) {AZ\\1};
\node[az] (az2) at (0,   0) {AZ\\2};
\node[az] (az3) at (1.6, 0) {AZ\\3};

%% Low-latency private links
\draw[<->, thick, saacolor!60] (az1) -- (az2);
\draw[<->, thick, saacolor!60] (az2) -- (az3);

%% Edge location to the right
\node[edge] (pop) at (4.8, 0) {Edge\\Location};
\node[lbl]  at (4.8, 0.65) {\tiny CloudFront / Route~53};

%% Internet arrow
\draw[->, thick, gray] (2.8, 0) -- node[above, font=\tiny]{internet} (pop.west);

\end{tikzpicture}
\end{center}

\begin{sapdepth}
\textbf{Multi-Region patterns (SAP focus):}
\begin{itemize}
  \item \textbf{Active--Active}: traffic served from multiple regions simultaneously (Route~53
        latency-based routing). Requires data synchronisation (DynamoDB Global Tables,
        Aurora Global write forwarding).
  \item \textbf{Active--Passive}: primary region handles all traffic; secondary region is
        on standby (warm standby or pilot light). Promotes on failover.
  \item \textbf{Data residency}: use AWS Organizations SCPs to \texttt{Deny} API calls
        outside approved regions — enforces sovereignty requirements.
  \item \textbf{AWS GovCloud}: isolated US regions for regulated workloads (FedRAMP, ITAR).
        Separate account required; accessed via specific endpoints.
\end{itemize}
\end{sapdepth}
